\section{Related Work}
\label{sec:related_work}

\subsection{Inherently Interpretable Tabular Models}
\label{subsec:related_interpretable}

Inherently interpretable (glass-box) models expose their decision logic without a
separate post-hoc explainer~\cite{rudin2019stop}; reviews of the tabular setting
\cite{obrien2024explainable} and of how such explanations are evaluated
\cite{nauta2023evaluating} give broader context.

\textbf{Linear and generalized linear models.} Linear and logistic regression are
among the oldest interpretable models, and each coefficient reads as the marginal
effect of a feature on the log-odds. They cannot represent non-linearities or
interactions without manual feature engineering.

\textbf{Decision trees and rule sets.} Trees and rule-based systems offer intuitive
\emph{IF--THEN} rules~\cite{wang2017bayesian,friedman2008predictive}, but reaching
competitive accuracy requires deep trees that are hard to follow, and axis-aligned
splits handle diagonal or correlated boundaries poorly.

\textbf{Generalized additive models.} GAMs decompose predictions into univariate
shape functions, $g(E[Y]) = \beta_0 + \sum f_i(x_i)$. Explainable Boosting
Machines~\cite{lou2012intelligible,caruana2015intelligible} add pairwise
interactions by shallow tree boosting, and Neural Additive
Models~\cite{agarwal2021neural} model each shape function with a small network.
Their cost grows sharply with many correlated features, many classes, and
higher-order interactions.

\subsection{Neural Rule Representations and Localized Activations}
\label{subsec:related_neural_rules}

Prototype networks~\cite{chen2019looks} and Radial Basis Function (RBF) networks
activate hidden units through distance-based centroids. Standard RBFs use isotropic
Euclidean distances that decay smoothly, so every neuron stays slightly active
everywhere. Piecewise-linear, compact-support activations instead drop to exactly
zero outside a defined region. Localized triangular (tent) activations suit such
units, but they have rarely been combined with bias-free projections and structural
pruning to extract discrete rule sets for tabular data. Sparse tabular neural
networks pursue a related goal with different mechanisms~\cite{raieli2026escaping}.
\modelname{} uses this geometry so that each neuron isolates a \emph{convex slab},
an explicit interval condition rather than an unbounded response.

\subsection{Neural Network Pruning and Structural Sparsity}
\label{subsec:related_pruning}

Neural network pruning was originally developed to compress large models and reduce memory and latency on hardware accelerators~\cite{han2015learning}. The Lottery Ticket Hypothesis by Frankle and Carbin~\cite{frankle2018lottery} demonstrated that dense, randomly initialized neural networks contain sparse sub-networks (``winning tickets'') that can reach comparable test accuracy when trained in isolation.

In \modelname{}, we use Iterative Magnitude Pruning (IMP) not only to compress the model but as a \emph{semantic sparsification mechanism}. In our architecture, setting an input connection $W_{f,hj}$ to zero directly removes feature $j$ from rule $h$, which shortens the rule. Also, when all input connections to a neuron are zeroed, the whole rule becomes inactive. Unlike conventional pruning pipelines that discard pruned units heuristically, we give a formal statement---the \emph{exact bias shifting for pruned rules}---which guarantees that disconnected rules can be absorbed into the output bias without altering model outputs.

\subsection{Scoring Systems and Integer Scorecards}
\label{subsec:related_scorecards}

Point-based scoring systems (scorecards) have been used for decades in medicine (e.g., the CHADS$_2$ score for stroke risk~\cite{gage2001validation}) and finance (e.g., credit bureau scorecards~\cite{hlongwane2024transparency}). In a scorecard, integer points are assigned to specific conditions, and the user simply adds the points together to obtain a risk category.

Modern methods such as Supersparse Linear Integer Models (SLIM)~\cite{ustun2016supersparse} and RiskSLIM~\cite{ustun2017optimized}, which has been applied to recidivism prediction~\cite{zeng2017interpretable}, train scorecards as a discrete optimization problem using Mixed-Integer Non-Linear Programming (MINLP). These methods produce provably optimal and highly sparse scorecards, but solving MINLP is NP-hard. As a result, training becomes very slow on datasets with many samples, many features, or multi-class targets. \modelname{} avoids this bottleneck by first training the continuous model with standard gradient descent and iterative pruning, and then converting the resulting compact rules into an integer scorecard via rule-adaptive scaling. The resulting prediction agreement is high but empirical rather than exact, and it is measured on the test partition. The recourse tool is likewise geometric rather than causal, a distinction that the counterfactual literature emphasizes~\cite{verma2024counterfactual}.

\subsection{Post-Hoc Explanations vs. Inherently Interpretable Models}
\label{subsec:related_posthoc}

Post-hoc explainability methods, particularly SHAP~\cite{lundberg2017unified} and LIME~\cite{ribeiro2016why}, have become a common way to explain black-box models. These tools build local surrogate models around individual predictions to estimate feature attributions.

Despite their popularity, post-hoc methods have fundamental limitations in
high-stakes domains~\cite{rudin2019stop}. They explain an approximation of the
model rather than its decision logic, so a diverging surrogate can mislead.
Estimating Shapley values is computationally expensive, and sampling-based
perturbations generate out-of-distribution points that distort attributions; the
same concern motivates dedicated evaluation protocols~\cite{nauta2023evaluating}.
Finally, some explainers can be fooled by adversarial wrappers that change the
explanation without changing the prediction~\cite{slack2020fooling}.

These limitations motivate inherently interpretable models such as \modelname{},
whose explanations are derived analytically from the model parameters.

\subsection{Positioning \modelname{} in the Literature}
\label{subsec:related_positioning}

Localized activations, magnitude pruning and integer scorecards all exist in
isolation. \modelname{} combines them so that the combination, rather than any
single component, carries the interpretability claim: bias-free projections with
triangular activations yield bounded convex-slab rules; pruning is semantic rather
than merely compressive and is backed by exact bias shifting; post-training
quantization produces integer scorecards without discrete optimization; and
decision decomposition, feature influence and recourse are available in closed
form. Relative to linear scorecards and GAMs, the representation is piecewise
linear and oblique, which is what allows it to cover localized non-axis-aligned
regions; relative to tree ensembles and deep networks, it exposes each decision as
a small number of signed rule votes rather than a surrogate explanation.
